\section*{Problem 1.2}

\subsection*{Part (a): Finding Commuting Matrices}

We seek two $2 \times 2$ matrices $A$ and $B$ such that $A \neq B$, neither is diagonal, $A \neq cB$ for any scalar $c$, and $AB = BA$.

Consider:
\begin{equation}
A = \begin{bmatrix} 1 & 1 \\ 0 & 1 \end{bmatrix}, \quad B = \begin{bmatrix} 1 & 2 \\ 0 & 1 \end{bmatrix}
\end{equation}

First, verify that $AB = BA$:
\begin{equation}
AB = \begin{bmatrix} 1 & 1 \\ 0 & 1 \end{bmatrix} \begin{bmatrix} 1 & 2 \\ 0 & 1 \end{bmatrix} = \begin{bmatrix} 1 & 3 \\ 0 & 1 \end{bmatrix}
\end{equation}

\begin{equation}
BA = \begin{bmatrix} 1 & 2 \\ 0 & 1 \end{bmatrix} \begin{bmatrix} 1 & 1 \\ 0 & 1 \end{bmatrix} = \begin{bmatrix} 1 & 3 \\ 0 & 1 \end{bmatrix}
\end{equation}

Therefore, $AB = BA$. Clearly, $A \neq B$, neither is diagonal, and $A \neq cB$ for any scalar $c$ (since the $(1,2)$ entries differ by a factor of 2, but the diagonal entries are equal).

\subsection*{Part (b): Finding Eigenvectors}

For matrix $A$, we solve $(A - \lambda I)v = 0$:
\begin{equation}
\det(A - \lambda I) = \det \begin{bmatrix} 1-\lambda & 1 \\ 0 & 1-\lambda \end{bmatrix} = (1-\lambda)^2 = 0
\end{equation}

Thus, $\lambda = 1$ is a repeated eigenvalue. For $\lambda = 1$:
\begin{equation}
(A - I)v = \begin{bmatrix} 0 & 1 \\ 0 & 0 \end{bmatrix} \begin{bmatrix} v_1 \\ v_2 \end{bmatrix} = \begin{bmatrix} 0 \\ 0 \end{bmatrix}
\end{equation}

This gives $v_2 = 0$, so the eigenvector is:
\begin{equation}
v_A = \begin{bmatrix} 1 \\ 0 \end{bmatrix}
\end{equation}

For matrix $B$, we solve $(B - \lambda I)v = 0$:
\begin{equation}
\det(B - \lambda I) = \det \begin{bmatrix} 1-\lambda & 2 \\ 0 & 1-\lambda \end{bmatrix} = (1-\lambda)^2 = 0
\end{equation}

Again, $\lambda = 1$ is a repeated eigenvalue. For $\lambda = 1$:
\begin{equation}
(B - I)v = \begin{bmatrix} 0 & 2 \\ 0 & 0 \end{bmatrix} \begin{bmatrix} v_1 \\ v_2 \end{bmatrix} = \begin{bmatrix} 0 \\ 0 \end{bmatrix}
\end{equation}

This gives $v_2 = 0$, so the eigenvector is:
\begin{equation}
v_B = \begin{bmatrix} 1 \\ 0 \end{bmatrix}
\end{equation}

\subsection*{Part (c): Shared Eigenvector}

As observed, both $A$ and $B$ share the eigenvector:
\begin{equation}
v = \begin{bmatrix} 1 \\ 0 \end{bmatrix}
\end{equation}

This demonstrates the fundamental result that commuting matrices share at least one eigenvector.

\section*{Problem 1.3}

We prove the three identities of Equation (1.26).

\subsection*{Identity 1: $(AB)^T = B^T A^T$}

Let $A$ be an $m \times n$ matrix and $B$ be an $n \times p$ matrix. Then $AB$ is $m \times p$, and $(AB)^T$ is $p \times m$. Let $C = AB$.

By definition of matrix multiplication and transpose:
\begin{equation}
C_{ij} = \sum_{k=1}^{n} A_{ik} B_{kj}
\end{equation}

\begin{equation}
(AB)^T_{ij} = C_{ji} = \sum_{k=1}^{n} A_{jk} B_{ki}
\end{equation}

Now consider $B^T A^T$. Since $B^T$ is $p \times n$ and $A^T$ is $n \times m$, their product is $p \times m$:
\begin{equation}
(B^T A^T)_{ij} = \sum_{k=1}^{n} B^T_{ik} A^T_{kj} = \sum_{k=1}^{n} B_{ki} A_{jk} = \sum_{k=1}^{n} A_{jk} B_{ki}
\end{equation}

Comparing equations (12) and (13), we have $(AB)^T_{ij} = (B^T A^T)_{ij}$ for all $i,j$. Therefore:
\begin{equation}
(AB)^T = B^T A^T
\end{equation}

\subsection*{Identity 2: $(AB)^{-1} = B^{-1} A^{-1}$}

Assume $A$ and $B$ are invertible square matrices of the same dimension. We must show that $B^{-1} A^{-1}$ is the inverse of $AB$.

By definition, we need to verify that $(AB)(B^{-1} A^{-1}) = I$ and $(B^{-1} A^{-1})(AB) = I$.

For the first product:
\begin{equation}
(AB)(B^{-1} A^{-1}) = A(BB^{-1})A^{-1} = A I A^{-1} = AA^{-1} = I
\end{equation}

For the second product:
\begin{equation}
(B^{-1} A^{-1})(AB) = B^{-1}(A^{-1}A)B = B^{-1} I B = B^{-1}B = I
\end{equation}

Since both products equal the identity matrix, we conclude:
\begin{equation}
(AB)^{-1} = B^{-1} A^{-1}
\end{equation}

\subsection*{Identity 3: $\text{Tr}(AB) = \text{Tr}(BA)$}

Let $A$ be an $m \times n$ matrix and $B$ be an $n \times m$ matrix. Then $AB$ is $m \times m$ and $BA$ is $n \times n$.

The trace of $AB$ is:
\begin{equation}
\text{Tr}(AB) = \sum_{i=1}^{m} (AB)_{ii} = \sum_{i=1}^{m} \sum_{j=1}^{n} A_{ij} B_{ji}
\end{equation}

The trace of $BA$ is:
\begin{equation}
\text{Tr}(BA) = \sum_{j=1}^{n} (BA)_{jj} = \sum_{j=1}^{n} \sum_{i=1}^{m} B_{ji} A_{ij}
\end{equation}

Since multiplication of scalars is commutative and the order of summation can be exchanged:
\begin{equation}
\text{Tr}(BA) = \sum_{j=1}^{n} \sum_{i=1}^{m} B_{ji} A_{ij} = \sum_{i=1}^{m} \sum_{j=1}^{n} A_{ij} B_{ji} = \text{Tr}(AB)
\end{equation}

Therefore:
\begin{equation}
\text{Tr}(AB) = \text{Tr}(BA)
\end{equation}

\section*{Problem 1.4}

We find the partial derivative of $\text{tr}(AB)$ with respect to $A$, where $A$ and $B$ are $m \times n$ and $n \times m$ matrices, respectively.

The trace of $AB$ is:
\begin{equation}
\text{tr}(AB) = \sum_{i=1}^{m} (AB)_{ii} = \sum_{i=1}^{m} \sum_{j=1}^{n} A_{ij} B_{ji}
\end{equation}

To find the derivative with respect to the matrix $A$, we compute the derivative with respect to each element $A_{kl}$:
\begin{equation}
\frac{\partial \text{tr}(AB)}{\partial A_{kl}} = \frac{\partial}{\partial A_{kl}} \sum_{i=1}^{m} \sum_{j=1}^{n} A_{ij} B_{ji}
\end{equation}

Since only the term with $i=k$ and $j=l$ contains $A_{kl}$:
\begin{equation}
\frac{\partial \text{tr}(AB)}{\partial A_{kl}} = B_{lk}
\end{equation}

Therefore, the derivative is:
\begin{equation}
\frac{\partial \text{tr}(AB)}{\partial A} = B^T
\end{equation}

This result states that the gradient of $\text{tr}(AB)$ with respect to $A$ is the transpose of $B$.

We can verify this using the cyclic property of the trace from Problem 1.3:
\begin{equation}
\text{tr}(AB) = \text{tr}(BA)
\end{equation}

Writing $\text{tr}(BA)$ in index notation:
\begin{equation}
\text{tr}(BA) = \sum_{j=1}^{n} \sum_{i=1}^{m} B_{ji} A_{ij} = \sum_{i=1}^{m} \sum_{j=1}^{n} B_{ji} A_{ij}
\end{equation}

Taking the derivative with respect to $A_{kl}$:
\begin{equation}
\frac{\partial \text{tr}(BA)}{\partial A_{kl}} = B_{lk} = (B^T)_{kl}
\end{equation}

This confirms our result:
\begin{equation}
\frac{\partial \text{tr}(AB)}{\partial A} = B^T
\end{equation}
